% Results-pending abstract: investigation language is intentional.
% Add measured findings only after completing the corresponding experiments.
\begin{abstract}
Multi-teacher on-policy distillation (MOPD) aims to consolidate specialized
capabilities in a single student. We examine three aspects of its
optimization: token balancing, gradient sparsity, and supervision
density. Starting from the MOPD learning objective, we first show how
loss averaging determines both domain weights and the relative weight
of long and short responses. A finite-batch gradient decomposition
separates these effects and motivates matched experiments on capability
balance. Second, we investigate whether different teachers concentrate
their gradients on overlapping parameter subsets within the same MOPD
student, and whether more different teacher distributions correspond to
more different subsets. Prior observations of sparse cumulative OPD
parameter changes motivate this question but do not establish gradient
sparsity; we measure gradients, optimizer steps, and cumulative changes
separately. Third, we compare sampled-token policy gradients with
teacher top-$k$ supervision in single-teacher and joint training,
using full-vocabulary reverse-KL gradients at shared prefixes as a
local reference. The planned experiments pair gradient concentration
with per-domain capability and account for sampling variability,
training exposure, and compute. Each question is developed together
with its experimental design, without presuming that sparse gradients,
low teacher overlap, or sampled supervision improve capability
integration. Empirical outcomes remain to be measured.
\end{abstract}
